% ================================================================
% ĐỀ THI 3: GIẢI TÍCH HÀM
% Mã học phần: 801047 | Thời gian: 90 phút
% ================================================================

\section*{Đề 3: Giải tích hàm}
\addcontentsline{toc}{section}{Đề 3: Giải tích hàm}

% ----------------------------------------------------------------
% HEADER THÔNG TIN ĐỀ THI
% ----------------------------------------------------------------
\noindent
\begin{minipage}[t]{0.45\textwidth}
    \begin{center}
        \textbf{TRƯỜNG ĐẠI HỌC SÀI GÒN}\\[0.3cm]
        \textbf{ĐỀ THI SỐ 3}
    \end{center}
\end{minipage}
\hfill
\begin{minipage}[t]{0.5\textwidth}
    \begin{center}
        \textbf{ĐỀ THI KẾT THÚC HỌC PHẦN}\\[0.2cm]
        Học phần: \textbf{Giải tích hàm}\\
        \textbf{Mã học phần:} 801047\\[0.2cm]
        \textit{Thời gian làm bài: 90 phút, không kể thời gian phát đề}
    \end{center}
\end{minipage}

\vspace{0.8cm}

% ----------------------------------------------------------------
% THÔNG TIN SINH VIÊN
% ----------------------------------------------------------------
\noindent
\rule{\textwidth}{0.5pt}\\[0.2cm]
Học kỳ: II \dotfill Năm học: 2025--2026 \dotfill\\[0.2cm]
Trình độ đào tạo: Đại học \dotfill Hình thức đào tạo: Chính quy \dotfill\\[0.2cm]
Họ tên sinh viên: \dotfill Mã số sinh viên: \dotfill\\[0.2cm]
\rule{\textwidth}{0.5pt}

\vspace{0.8cm}

% ----------------------------------------------------------------
% LƯU Ý CHO SINH VIÊN
% ----------------------------------------------------------------
\begin{center}
    \fbox{%
        \parbox{0.8\textwidth}{%
            \centering
            \textbf{LƯU Ý:} Sinh viên \textbf{KHÔNG} sử dụng tài liệu.\\
            Cán bộ coi thi không giải thích gì thêm.
        }%
    }
\end{center}

\vspace{1cm}

% ================================================================
% NỘI DUNG ĐỀ THI
% ================================================================

% ----------------------------------------------------------------
% CÂU 1 (2.0 ĐIỂM)
% ----------------------------------------------------------------
\noindent\textbf{\large Câu 1. (2,0 điểm)}

\vspace{0.3cm}

\textbf{a)} Trên không gian véctơ $\mathbb{R}^2$, xét hàm số:
\[
    \|(x, y)\| = 2|x| + |y| \quad \text{với mọi } (x, y) \in \mathbb{R}^2.
\]
Chứng minh rằng hàm số này là một chuẩn trên $\mathbb{R}^2$.

\vspace{0.4cm}

\textbf{b)} Hãy mô tả hình học quả cầu đơn vị đóng:
\[
    \overline{B}(0, 1) = \{ (x, y) \in \mathbb{R}^2 : \|(x, y)\| \leq 1 \}
\]
trong mặt phẳng tọa độ $Oxy$.

% ----------------------------------------------------------------
% CÂU 2 (2.0 ĐIỂM)
% ----------------------------------------------------------------
\vspace{0.6cm}
\noindent\textbf{\large Câu 2. (2,0 điểm)}

\vspace{0.3cm}

Xét không gian các hàm số liên tục $X = C([0, 1])$ với chuẩn:
\[
    \|f\|_\infty = \max_{x \in [0, 1]} |f(x)|.
\]

Xét ánh xạ $T : X \to X$ xác định bởi:
\[
    (Tf)(x) = x \int_0^1 y^2 f(y) \, dy, \quad \text{với mọi } f \in X.
\]

\textbf{a)} Chứng minh $T$ là một toán tử tuyến tính liên tục.

\vspace{0.3cm}

\textbf{b)} Tính chuẩn toán tử $\|T\|_{B(X, X)}$.

% ----------------------------------------------------------------
% CÂU 3 (3.0 ĐIỂM)
% ----------------------------------------------------------------
\vspace{0.6cm}
\noindent\textbf{\large Câu 3. (3,0 điểm)}

\vspace{0.3cm}

Xét không gian $\mathbb{R}^4$ với chuẩn Euclide:
\[
    \|(x_1, x_2, x_3, x_4)\| = \sqrt{x_1^2 + x_2^2 + x_3^2 + x_4^2}.
\]

Xét không gian con:
\[
    M = \{(x_1, x_2, x_3, x_4) \in \mathbb{R}^4 : x_1 = x_2, \; x_2 = 2x_3, \; x_3 = -x_4\}
\]

và ánh xạ $f : M \to \mathbb{R}$ cho bởi:
\[
    f(x_1, x_2, x_3, x_4) = x_1 + x_2 - x_3 + 2x_4, \quad \forall x \in M.
\]

\textbf{a)} Chứng minh $f$ là phiếm hàm tuyến tính liên tục trên $M$.

\vspace{0.3cm}

\textbf{b)} Chứng minh rằng tồn tại phiếm hàm tuyến tính liên tục $F : \mathbb{R}^4 \to \mathbb{R}$ sao cho:
\[
    F|_M = f \quad \text{và} \quad \|F\|_{B(\mathbb{R}^4, \mathbb{R})} = \|f\|_{B(M, \mathbb{R})}.
\]
Tìm công thức cụ thể của toán tử $F$.

% ----------------------------------------------------------------
% CÂU 4 (3.0 ĐIỂM)
% ----------------------------------------------------------------
\vspace{0.6cm}
\noindent\textbf{\large Câu 4. (3,0 điểm)}

\vspace{0.3cm}

Xét không gian Hilbert $H = L^2([-1, 1])$ với tích vô hướng:
\[
    \langle f,g \rangle = \int_{-1}^1 f(x)g(x)\,dx.
\]

Gọi $M_e$ và $M_o$ lần lượt là tập hợp các hàm số chẵn và hàm số lẻ thuộc $H$.

\vspace{0.3cm}

\textbf{a)} Chứng minh rằng $M_e$ và $M_o$ là hai không gian con trực giao của $H$.

\vspace{0.3cm}

\textbf{b)} Xét hàm số $f(x) = x^3 + 3x^2 - 2x + 1$. Hãy tìm hình chiếu trực giao của $f$ lên không gian con $M_o$ và tính khoảng cách từ $f$ đến $M_o$.

% ----------------------------------------------------------------
% KẾT THÚC ĐỀ THI
% ----------------------------------------------------------------
\vspace{1.5cm}
\begin{center}
    \rule{0.5\textwidth}{0.5pt}\\[0.3cm]
    \textbf{\large --- HẾT ---}\\[0.3cm]
    \rule{0.5\textwidth}{0.5pt}
\end{center}

\newpage
